%\documentclass[11pt,a4paper,oneside]{report}             % Single-side
\documentclass[11pt,a4paper,twoside,openright]{report}  % Duplex

%\PassOptionsToPackage{chapternumber=Huordinal}{magyar.ldf}
\usepackage[T1]{fontenc}
\usepackage[utf8]{inputenc}
\usepackage{times}
\usepackage{amsmath}
\usepackage{amssymb}
\usepackage{enumerate}
\usepackage[thmmarks]{ntheorem}
\usepackage{graphics}
\usepackage{epsfig}
%\usepackage{listings}
\usepackage{listingsutf8} % Ezzel a csomaggal működik rendesen a listings UTF-8 esetén is.
\usepackage{color}
%\usepackage{fancyhdr}
\usepackage{lastpage}
\usepackage{anysize}
\usepackage[english,magyar]{babel}
\usepackage{sectsty}
\usepackage{setspace}  % Ettol a tablazatok, abrak, labjegyzetek maradnak 1-es sorkozzel!
\usepackage[hang]{caption}
\usepackage{hyperref}

% Táblázatoknak:
\usepackage{colortbl}
\usepackage{longtable}

\usepackage{verbatim}

% Egymás melleti képekhez:
\usepackage{subfig}

%--------------------------------------------------------------------------------------
% Main variables
%--------------------------------------------------------------------------------------
\newcommand{\vikszerzo}{Nádudvari György}
\newcommand{\vikkonzulens}{Huszerl Gábor}
\newcommand{\vikcim}{Oktatástámogató rendszerek kiszolgáló infrastruktúrájának felügyeleti lehetőségei}
\newcommand{\viktanszek}{Méréstechnika és Információs Rendszerek Tanszék}
\newcommand{\vikdoktipus}{Szakdolgozat}

%--------------------------------------------------------------------------------------
% Page layout setup
%--------------------------------------------------------------------------------------
% we need to redefine the pagestyle plain
% another possibility is to use the body of this command without \fancypagestyle
% and use \pagestyle{fancy} but in that case the special pages
% (like the ToC, the References, and the Chapter pages)remain in plane style

\pagestyle{plain}
%\setlength{\parindent}{0pt} % áttekinthetőbb, angol nyelvű dokumentumokban jellemző
%\setlength{\parskip}{8pt plus 3pt minus 3pt} % áttekinthetőbb, angol nyelvű dokumentumokban jellemző
\setlength{\parindent}{12pt} % magyar nyelvű dokumentumokban jellemző
\setlength{\parskip}{0pt}    % magyar nyelvű dokumentumokban jellemző

\marginsize{35mm}{25mm}{15mm}{15mm} % anysize package
\setcounter{secnumdepth}{0}
\sectionfont{\large\upshape\bfseries}
\setcounter{secnumdepth}{2}
\singlespacing
\frenchspacing

%--------------------------------------------------------------------------------------
%	Setup hyperref package
%--------------------------------------------------------------------------------------
\hypersetup{
    bookmarks=true,            % show bookmarks bar?
    unicode=false,             % non-Latin characters in Acrobat’s bookmarks
    pdftitle={\vikcim},        % title
    pdfauthor={\vikszerzo},    % author
    pdfsubject={\vikdoktipus}, % subject of the document
    pdfcreator={\vikszerzo},   % creator of the document
    pdfproducer={Producer},    % producer of the document
    pdfkeywords={keywords},    % list of keywords
    pdfnewwindow=true,         % links in new window
    colorlinks=true,           % false: boxed links; true: colored links
    linkcolor=black,           % color of internal links
    citecolor=black,           % color of links to bibliography
    filecolor=black,           % color of file links
    urlcolor=black             % color of external links
}

%--------------------------------------------------------------------------------------
% Set up listings
%--------------------------------------------------------------------------------------
\lstset{
	basicstyle=\scriptsize\ttfamily, % print whole listing small
	keywordstyle=\color{black}\bfseries\underbar, % underlined bold black keywords
	identifierstyle=, 					% nothing happens
	commentstyle=\color{white}, % white comments
	stringstyle=\scriptsize\sffamily, 			% typewriter type for strings
	showstringspaces=false,     % no special string spaces
	aboveskip=3pt,
	belowskip=3pt,
	columns=fixed,
	backgroundcolor=\color{lightgray},
	extendedchars=\true, % Kiegészítés az UTF-8 támogatáshoz
    inputencoding=utf8,
} 		
\def\lstlistingname{lista}	

%--------------------------------------------------------------------------------------
%	Some new commands and declarations
%--------------------------------------------------------------------------------------
\newcommand{\code}[1]{{\upshape\ttfamily\scriptsize\indent #1}}

% define references
\newcommand{\figref}[1]{\ref{fig:#1}.}
\renewcommand{\eqref}[1]{(\ref{eq:#1})}
\newcommand{\listref}[1]{\ref{listing:#1}.}
\newcommand{\sectref}[1]{\ref{sect:#1}}
\newcommand{\tabref}[1]{\ref{tab:#1}.}

\DeclareMathOperator*{\argmax}{arg\,max}
%\DeclareMathOperator*[1]{\floor}{arg\,max}
\DeclareMathOperator{\sign}{sgn}
\DeclareMathOperator{\rot}{rot}
\definecolor{lightgray}{rgb}{0.95,0.95,0.95}

\author{\vikszerzo}
\title{\viktitle}

\begin{comment}
\includeonly{
	guideline,%
	project,%
	titlepage,%
	declaration,%
	abstract,%
	introduction,%
	fejezet_lmsek,
	fejezet_lmsek_it_struk,
	acknowledgement,%
	appendices,%
}
\end{comment}
%--------------------------------------------------------------------------------------
%	Setup captions
%--------------------------------------------------------------------------------------
\captionsetup[figure]{
%labelsep=none,
%font={footnotesize,it},
%justification=justified,
width=.75\textwidth,
aboveskip=10pt}

\renewcommand{\captionlabelfont}{\small\bf}
\renewcommand{\captionfont}{\footnotesize\it}

%###########################################
% Saját eszközök:
%
\definecolor{todobgszin}{rgb}{0.64,0.78,0.22}
\definecolor{todofrszin}{rgb}{0.00,0.50,0.00}
%\definecolor{MyTableColor}{rgb}{0.38,0.28,0.25}
%\definecolor{MyTableColor}{rgb}{0.74,0.85,0.34} % June Bud
\definecolor{MyTableColor}{rgb}{0.82,0.89,0.19} % Pear

\newcommand{\todo}[1]{\fcolorbox{todofrszin}{todobgszin}{\emph{TODO: #1}}}
\newcommand{\angolul}[1]{\foreignlanguage{english}{#1}}

\newenvironment{sajat_itemize}
{
	\begin{itemize}
	\setlength{\itemsep}{0pt}
}
{
	\end{itemize}
}

%--------------------------------------------------------------------------------------
% Table of contents and the main text
%--------------------------------------------------------------------------------------
\begin{document}
\singlespacing
%\include{project}

\pagenumbering{arabic}
\onehalfspacing

\include{titlepage}
\tableofcontents\vfill
\include{declaration}
%----------------------------------------------------------------------------
% Abstract in hungarian
%----------------------------------------------------------------------------
\chapter*{Kivonat}\addcontentsline{toc}{chapter}{Kivonat}

Napjainkban egyre több területen alkalmazzák sikerrel az elektronikus oktatást. Ezen rendszerek üzemeltetése nagyrészt a szokásos -- egyéb rendszerekre is jellemző -- feladatokat adja a kiszolgáló infrastruktúrát felügyelők számára, azonban az erőforrásigényük ingadozása igen speciális. Ezért érdemes áttekinteni, hogy az oktatástámogató rendszerek használatához kapcsolódóan keletkező információk milyen segítséget jelenthetnek az ezeket a rendszereket kiszolgáló infrastruktúra felügyeleténél.

A téma nehézségét adja, hogy nem rendelkezünk az oktatást támogató rendszerek erőforrásigényei és az IT infrastruktúra elérhető erőforrásai közötti megfeleltetésekkel.

Feladatom volt, hogy megismerkedjek ezen oktatástámogató rendszerek használati területeivel, felépítésével, erőforrásigényeivel és azok várható változásaival. Megvizsgáltam a felhőalapú infrastruktúra erőforrás-kezelési lehetőségeit, és megpróbáltam proaktív megoldást adni az oktatástámogató rendszerek erőforrás-kezelésére.
\vfill

%----------------------------------------------------------------------------
% Abstract in english
%----------------------------------------------------------------------------
\chapter*{Abstract}\addcontentsline{toc}{chapter}{Abstract}

Nowadays the e-learning is successfully applied in many field. The operating of these systems mostly provides the usual tasks for the managers of the serving infrastucture, however, there is the fluctuation of the resource needs which is very special. It's worth to reviewing that, how these generated information from the use of learning management systems can help the management of the serving infrastructure.

In this subject the difficulty is that, we do not have matching between the resource needs of an e-learning system and the available resources of the IT infrastructure.

My task was to familiarize myself with the support areas, structure, resource needs of these learning management systems and the changes of these resource needs. I examined the resource management capabilities of cloud-based infrastructure, and tried to find a proactive solution to supporting the e-learning systems of resource management.

\vfill


%----------------------------------------------------------------------------
\chapter*{Bevezető}\addcontentsline{toc}{chapter}{Bevezető}
%----------------------------------------------------------------------------

Napjainkban egyre több területen jelenik meg az elektronikus oktatás. Manapság már nem csak a közoktatásban, de az iparban is szívesen alkalmazzák a számítógép és Internet által támogatott tanulást, továbbképzést. Ezzel együtt növekszik az igény az ilyen rendszerek megfelelő, nagy rendelkezésre állást biztosító üzemeltetése iránt.

Az elektronikus oktatástámogató alkalmazások erőforrásigényének változása speciális tulajdonságokkal rendelkezik. Egyrészről ezen ingadozások egy része determinisztikus, nagy pontossággal előre meghatározható, másrészről viszont az üzemeltetés oldaláról igényt tartanak az állandó megfigyelésre, analizálásra és időbeni reakciókra.

Az oktatási rendszerek nagyon gyorsan terjednek, ugyanakkor még nincs érdemi múltjuk, ezen felül olyan új technológiákat alkalmazhatnak, mint pl. a felhőalapú infrastruktúrák, ezért ezen a területen még nem rendelkezünk megfelelő információkkal arról, hogy hogyan tudjuk az erőforrás-kihasználtsági problémákat effektíven kezelni. Ezeket a problémákat és lehetséges megoldásaikat szeretném megmutatni írásomban.

Szakdolgozatomban először szeretném bemutatni az ún. tanulásmenedzsment vagy oktatástámogató rendszereket, azok főbb tulajdonságait. Kitérek a piacon található kész rendszerekre, azok megvalósítási platformjaira.

A második részben az ismert három rétegű architektúra példáján kiindulva bemutatom ezen rendszerek üzemeltetési IT infrastruktúráját, majd egy példa rendszer felépítését.

A harmadik fejezetben az oktatástámogató rendszerek erőforrásigényeit, azok változásait és a változások lehetséges okait vizsgálom. Ugyanebben a részben megmutatom ezen igényváltozások modellezési lehetőségeit.

A következő fejezetben a klasszikus IT infrastruktúra megbízhatóságának, a rendelkezésre állás növelésének lehetőségeit járom körül, majd bemutatom az ún. felhő alapú megoldásokat és azok jellemzőit, szolgáltatásait.

Az utolsó részben a rendszer lehetséges menedzselési típusait mutatom be, részletesebben megvizsgálva a proaktív menedzsment jellemzőit, lehetséges alkalmazásait az oktatástámogató rendszerek üzemeltetése esetén.
 
\chapter{Tanulásmenedzsment rendszerek}
LMS-nek (Learning Management System) nevezzük azokat a szoftver alkalmazásokat, amelyek automatizálják az oktatás adminisztrációját, követését, az online kurzusok és az azokkal kapcsolatos események, anyagok kezelését.

Egy robusztus LMS-nek képesnek kell lennie \cite{link:ell}:
\begin{itemize}
\item központosított és automatizált adminisztrációra
\item önkiszolgáló és önálló irányítású szolgáltatások nyújtására
\item oktatási anyagok gyors összeállítására és szállítására
\item konszolidált képzési kezdeményezésekre skálázható, web alapú platformon
\item a portabilitás és a szabványok támogatására
\item személyre szabott tartalom előállítására és a tudás újrafelhasználásának lehetővé tételére.
\end{itemize}

\Aref{tab:openlms}.~táblázatban néhány ismertebb nyílt forrású LMS-t soroltam fel \cite{link:lms}. Érdemes megnézni, hogy általában PHP, MySQL technológiák segítségével valósítják meg ezeket a rendszereket.

\definecolor{MyTableColor}{rgb}{0.38,0.28,0.25} 

\begin{table}[h]
	\caption{Néhány elterjedtebb LMS}
	\centering
	\begin{tabular}{| p{3cm} | p{4cm} | p{2cm} | p{4cm} |}
		\hline
		\rowcolor{MyTableColor} \textbf{Név} & \textbf{Projekt oldal} & \textbf{Használt prog.~nyelv} & Támogatott adatbázisok \\
		\hline
		Sakai & \href{http://www.sakaiproject.org/}{http://www.sakaiproject.org/} & Java & MySQL,~Oracle,~DB2 \\
		\hline
		Moodle & \href{http://moodle.org/}{http://moodle.org/} & PHP & MySQL, PostgreSQL, MSSQL, Oracle \\
		\hline
		OLAT &  \href{http://www.olat.org/}{http://www.olat.org/} & Java & MySQL,~PostgreSQL \\
		\hline
		Instructure - Canvas & \href{http://www.instructure.com/}{http://www.instructure.com/} & Ruby on Rails & PostgreSQL \\
		\hline
		ILIAS & \href{http://www.ilias.de/docu/}{http://www.ilias.de/docu/} & PHP & MySQL (Oracle a 4.2-es verzióban lesz) \\
		\hline
		Dokeos & \href{http://www.dokeos.com/}{http://www.dokeos.com/} & PHP & MySQL \\
		\hline
		Chamilo & \href{http://www.chamilo.org/}{http://www.chamilo.org/} & PHP & MySQL \\
		\hline		
		Claroline & \href{http://www.claroline.net/}{http://www.claroline.net/} & PHP & MySQL \\
		\hline
	\end{tabular}
	\label{tab:openlms}
\end{table}

\chapter{Tanulásmenedzsment rendszerek IT infrastruktúrája}
\section{A három rétegű architektúra}
A webes LSM-ek általában a három rétegű architektúrát követik. Ez a három réteg a webkiszolgáló, az adatbázis és az alkalmazás réteg. A réteges szerkezetnek köszönhetően ezek a rendszerek jól skálázhatóak, hiszen az egyes rétegekben megtalálható szolgáltatásokat gyakran ilyen funkciókkal alakítják ki.
A rétegeknek nem szükséges fizikailag is külön hardverre kerülni (sőt az alkalmazás- és webkiszolgáló-réteg esetében ez  nem is mindig lehetséges), így a legegyszerűbb kialakítás akár egy számítógépet is igénybe vehet. Ez a megoldás egy viszonylag erős konfiguráció és alacsony felhasználószám esetén működhet.
\subsection{A webkiszolgáló-réteg}
\todo{Webszerver}
\subsection{Az alkalmazásréteg}
\todo{Alkalmazás szerver}
\subsection{Az adatbázisréteg}
\todo{DB}

És itt elő is kerül az elterjedtebb LMS-ek előnye: a könnyű telepítés, beüzemelés. Vegyük például

\section{Egy példa}
\subsection{A Moodle rendszer}
Az e-learning rendszerek közül a legelterjedtebb a Moodle (Modular Object-Oriented Dynamic Learning Environment) (\href{http://moodle.org}{http://moodle.org}) nevű tanulásmenedzsment rendszer. Önálló laboratóriumi munkám során ennek a rendszernek a részletesebb megismerése volt az egyik cél.

A Moodle  egy ingyenes és nyílt forrású LMS vagy VLE (Virtual Learning Environment, Virtuális oktató környezet). 2010 októberében 49952 regisztrált Moodle oldal létezett, amelyek összesen mintegy 37 millió felhasználót szolgáltak ki. A legnagyobb rendszertelepítések közé tartozik a tajvani Ming Chuan Egyetem több, mint 63.000 regisztrált, maximálisan 33.000 bejelentkezett felhasználóval naponta.

\todo{Forrás http://docs.moodle.org/20/en/Installations\_30000\_plus}

Maga a rendszer tervezéséből és implementálásából eredően portábilis, hála a PHP nyelvnek. Módosítás nélkül telepíthető Unix, Linux variánsokra, FreeBSD-re, Windows-ra, Mac OS X-re, NetWare-re és egyéb rendszerekre, amelyek támogatják a PHP-t és az ismertebb adatbázis-kezelő rendszereket.

A Moodle architektúrájában elválik a natív adatbázis-kezelő a felsőbb rétegektől. Közöttük egy ADOdb adatbázis absztrakciós réteg található. Az ADOdb egy különálló projekt (\href{http://adodb.sourceforge.net/}{http://adodb.sourceforge.net/}), amely a legelterjedtebb adatbázis-kezelőket támogatja. \Aref{fig:moodlearch}.~ábrán egy korábbi verzió architektúrája látható.

\begin{figure}[h!]
\centering
\includegraphics[width=0.8\textwidth]{figures/moodlearch.png}
\caption{A Moodle architektúrája \label{fig:moodlearch}}
\end{figure}

Az architekturális felépítésből látható, hogy a Moodle az ADOdb által nyújtott rétegen keresztül egységesen éri el az eltérő adatbázis megvalósításokat anélkül, hogy foglalkoznia kellene az ezekből eredő problémákkal.  

Az interoperabilitás több Moodle képességben megjelenik, ilyenek például:
\begin{itemize}
\item autentikáció LDAP-on, Shibboleth-en keresztül
\item kérdések/kérdéssorok importálása/exportálása több formátumban (pl. XML, XHTML stb.)
\item erőforrások kezelése (SCORM, AICC)
\item integráció más tartalom kezelő rendszerekkel (Drupal, Postnuke).
\end{itemize}

A Moodle fontos tulajdonsága a modularitás, vagyis a rendszer funkcionalitásának könnyű bővíthetőségek. Ezeket plug-inokkal valósítják meg, amelyek fejlesztését különböző API segítik. Ilyen plug-inok pl. a különböző jelentések (reportok), blokkok, portfóliók, tárolók (repository-k), kereső modulok. Ezek közül nagyon sok megtalálható az alap Moodle telepítésben is, de saját magunk is írhatunk hasonló kiegészítőket.

\subsection{A Moodle egy lehetséges infrastruktúra kialakítása}

\chapter{Tanulásmenedzsment rendszerek erőforrásigényei}

Minden informatikai rendszernek van erőforrásigénye, vagyis az a minimális hardver konfiguráció, amelyen a rendszer bizonyos számú kérést ki tud szolgálni adott időn belül. Az erőforrások mérete befolyásolja a kiszolgálható kérések számát. Ha a kérések száma növekszik, akkor növekedhet az erőforrásigény is. A kiszolgálandó kérések típusától, fajtájától függ, hogy pontosan milyen erőforrásokra van igény.
Egy informatikai rendszer esetében a leggyakrabban emlegetett, figyelt, használt erőforrások:
\begin{sajat_itemize}
\item CPU (pl. mennyi processzoridőt használ a rendszer, mennyit az alkalmazás),
\item memória (pl. mennyi memóriát használ az alkalmazásunk),
\item tárhely (pl. mennyi tárhelyet foglal el az alkalmazásunk az adataival együtt),
\item hálózat (pl. mennyire van kihasználva a sávszélesség, mennyi a ki és be irányú forgalom),
\item I/O sávszélesség
\item stb.
\end{sajat_itemize}

A következőekben egy alapvető rendszer és az arra ráépülő, a háromrétegű architektúrát megvalósító alkalmazások erőforrásigényeit vizsgálom meg, majd pedig ezen igények lehetséges változásait és azok okait gyűjtöm össze egy LMS rendszer esetében. 

\section{Alapvető igények}

Egy informatikai rendszer teljes erőforrásigénye sok részletből tevődhet össze. Alapvetően meghatározza az operációs rendszer igénye, az erre épülő szolgáltatások (adatbázis szerver, webszerver, levelezőszerver, egyebek) és a főszolgáltatás platformja.
A ma használatos szerver operációs rendszerek alapkövetelményének számít \cite{ws2008sr,ubuntuminhr,redhatcaplim}:

\begin{sajat_itemize}
\item min. 1-1,5 GHz-es CPU (x86 vagy x64), de 2 GHz az ajánlott,
\item min. 512 MB memória, de 2 GB ajánlott\footnote{Red Hat Linux esetében pl. 1 GB/processzor egység},
\item min. 10 GB tárkapacitás.
\end{sajat_itemize}

Gyakran ajánlják a szerverek operációs rendszerének a Linux rendszereket, mivel ezek általában kevesebb memóriát használnak fel, így több marad a többi szolgáltatás számára. 

Nem csak az LMS-ek, de minden webes alkalmazás esetén is különböző platformokat találunk az alkalmazásszerver oldalán. Az igen elterjedt, PHP mellett az üzleti életben gyakran Java alapokon fejlesztenek, de nem szabad megfeledkezni a .NET-es, Ruby on rails-es vagy a Python-os megvalósításokról sem. Az egyes platformok különböznek teljesítményükben, ''gyorsaságukban'', megbízhatóságukban és erőforrásigényükben is. Az egyik legelterjedtebb webes alkalmazás platform a PHP, mivel ez a nyelv szinte minden operációs rendszeren, szinte minden webszerverrel együtt tud működni, és emellett igen jól skálázható.

\subsection{Webszerverek erőforrásigénye}

Ahhoz, hogy el tudjunk indulni a webszerverek erőforrásigényeinek vizsgálatában, először meg kell értenünk, hogyan is működnek ezek az alkalmazások. Egy webszerver általában egy többfolyamatos (\angolul{multi-process}) vagy többszálas (\angolul{multi-threaded}) modell szerint működik\cite{Lu:2001:FCA:882481.883781}. Ezek a feldolgozó folyamatok vagy szálak igény esetén jönnek létre, vagy egy tárolóban előre létrehozott számban várják a beérkező TCP kapcsolatokat, hogy kiszolgálhassák azokat. A széleskörűen használt Apache webszerver a többfolyamatos, készletes modellt alkalmazza.

A HTTP 1.0 esetében minden egyes kérés egy TCP kapcsolatot jelentett, amely így nagyszámú konkurens kapcsolatot igényelt. Ezen javítva a HTTP 1.1-es verziójában megjelent a perzisztens kapcsolat, amely lehetővé teszi, hogy egy kapcsolatba több kérés is belekerüljön\cite{rfc1945,rfc2616}.

Ezek a perzisztens kapcsolatok egy új típusú szűk keresztmetszetet hoztak be a szerverekbe. Amióta a kiszolgáló folyamat egy perzisztens kapcsolathoz köthető, a CPU kihasználtsága nagyon alacsony. Ezen alacsony kihasználtságon a végrehajtó folyamatok számának növelésével segíthetünk, ám ekkor a virtuális memória kezdhet el vergődni (\angolul{thrashing}). Ezt csak viszonylag sok elérhető memóriával orvosolhatjuk. Ebből következik, hogy a webszerverek inkább \mbox{memória-,} mint processzorigényesek. 

\subsection{Adatbázisszerverek erőforrásigénye}

Az adatbázisok erőforrásigénye igen összetett, hiszen egyik részről a lekérdezéseket, tranzakciókat optimalizálni kell (CPU terhelés), az adatokat a háttértárról be kell olvasni, vagy oda ki kell írni (I/O terhelés), és az ezeket az igények csökkentő technikák memóriaigénye növekedhet.

Az adatbázisszerverek nagy része támogatja az ún. \angolul{connection pooling} megoldást, amely a webszerverek perzisztens kapcsolat-kezeléséhez hasonló, és ezzel együtt ebben az esetben is az erőforrásigény a memória felé tolódik el. Tehát összességében az adatbázisszerverek is inkább memória és lemez I/O igényesek.

Néhány adatbázisnál előnyként jelenik meg, hogy a \angolul{connection pooling}gal úgy tudjuk a teljesítményt növelni, hogy ahhoz nem kell az alkalmazásréteghez hozzányúlni. A PostgreSQL adatbázis-kezelőhöz fejlesztett pgpool (\href{http://www.pgpool.net}{http://www.pgpool.net}) például egy ilyen megoldás, amellyel akár kétszeres kiszolgálási értékeket is elérhetünk.

\section{Egy példarendszer: a Moodle}
A Moodle hardveres erőforrásigényei a következők:
\begin{sajat_itemize}
\item min. 160 MB tárolókapacitás (csak a rendszernek \footnote{A teljes tárolókapacitás függ az oktatási anyag mennyiségétől, méretétől, típusától (pl. oktatóvideók)}),
\item min. 256 MB memóriakapacitás (1 GB ajánlott).
\end{sajat_itemize}
Érdemes megjegyezni, hogy a hivatalos Moodle oldalon 1 GB memóriát ajánlanak 50 felhasználó egyszerre történő kiszolgálásához \cite{moodleinst}.
Természetesen ez még nem a teljes rendszer igénye, hiszen a webkiszolgálónak és az adatbázisszervernek is van erőforrásigénye, de egy viszonylag alacsony felhasználói számmal rendelkező szervernek elegendő az 1 GB memóriakapacitás.

A processzor oldaláról nincs különösebb megkötés, a manapság már általánosnak mondható kétmagos, 1,5 GHz-es processzor már elégnek bizonyulhat a rendszer működéséhez.

\section{Igények terhelésváltozások esetén}

Az LMS-ek esetében a terhelésváltozások okozói elsősorban maguk a felhasználók, hiszen ők azok, akik időszakosan megnövekedett kérésszámot intézhetnek a rendszerünk felé. Érdemes lehet ezeket az időszakokat és a hozzájuk kapcsolódó viselkedési változásokat összegyűjteni.

Egy LMS életében a következő erőforrásigényes időszakokat különböztethetjük meg:
\begin{sajat_itemize}
\item kurzus-/vizsgajelentkezési időszak,
\item kurzussal kapcsolatos feladatok beadási határideje,
\item kurzus online teszt, vagy vizsga kitöltés (határ)ideje,
\item egyéb a kurzussal kapcsolatos offline számonkérés,
\item online előadás közvetítés,
\item audiovizuális tananyagokkal rendelkező kurzus számonkérésének ideje, 
\end{sajat_itemize}
és még hosszasan sorolhatnánk. A következőekben megpróbálom összefoglalni, hogy ezek közül melyik milyen módon terheli meg a rendszerünk erőforrásait.

\subsubsection{Kurzus-/vizsgajelentkezési időszak\protect\footnote{Ezt az alfejezetet a BME KTH Neptun Üzemeltetés által kiadott tájékoztatók alapján írtam, amelyek a rendszer bejelentkező oldalán, a "Letölthető dokumentumok" résznél találhatóak.}}

Sajnos a BME hallgatóinak nem ismeretlen a Neptun Egységes Tanulmányi Rendszer vizsga- és tárgyjelentkezési időszakokban történő kiesése. Ezekben az időszakokban rohamosan megnő az egyszerre aktív felhasználók száma, amivel a rendszer nem tud megbirkózni. Az üzemeltetés mindent megtesz annak érdekében, hogy a rendszer működőképes legyen a több ezer felhasználó esetén is, de ha valamelyik rétegben egy-egy erőforrás nem tudja feldolgozni a kéréseket, akkor az az egész rendszerre kihatással van.

A rendszer terheléselosztását egy virtualizált gépre telepített szoftver végzi. A 2011. év januári tárgyfelvételi időszakában ez az ingyenes, de több helyen sikerrel alkalmazott szoftver 4000-5000 konkurens felhasználót még ki tudott szolgálni, de az efölötti kérések tömegével már nem volt képes megbirkózni.

A rendszer úgy van kialakítva, hogy a terheléselosztó építi ki a kapcsolatot a webszerverekkel, viszont amikor az összeomlik, a már felépített kapcsolatok is elvesznek. A felhasználók ebből annyit tapasztaltak, hogy a rendszer csak hosszú idő után engedte őket belépni a rendszerbe, majd ezután gyakran ki is dobta őket.

A felhasználók a negatív tapasztalatokra még agresszívebben próbálták elérni a rendszert (folyamatos oldallekérések, frissítések), amely viselkedés csak növelte a kiszolgálatlan kérések számát. Mindeközben a webszerverek és az adatbázisszerverek kihasználtsága a nullához közelített, hiszen a terheléselosztó nem volt képes a kéréseket továbbítani feléjük.

A Neptun rendszer a júniusi tárgyfelvételi időszakot is zökkenővel viselte. A június 7-én 18 órakor kezdődő felhasználói rohamot, ami 6000 egyidejű, aktív felhasználót jelentett, ezúttal a terheléselosztó jól viselte, ám pontosan emiatt az adatbázisszerver nem tudod megküzdeni a nagy számú kéréssel, és megtelt a memóriája.

Ennek következtében az üzemeltetésnek le kellett állítani az egész rendszert, újrakonfigurálni az adatbázisszervert (szerencsére mivel virtualizált rendszerről beszélünk, ezért viszonylag egyszerű megnövelni a memóriáját), majd újraindítani azt és a webszervereket. Ezután a rendszer már megfelelően működött, alig 10 perc leforgása alatt 6000 felhasználó tudott belépni, és 22500 tárgyat felvenni.

A 2011. augusztus végi tárgyfelvételi időszak jól jellemzi, hogy mennyire függ egy ilyen rendszer a felhasználói viselkedéstől.  Mivel a hallgatóknak már lehetőségük volt a nyár elején felvenni a következő féléves tárgyaikat, ezért a véglegesítési időszakban elmaradt a szokásos terhelés.

A 2011. november 28-án 18 órakor kezdődő vizsgajelentkezés során megint problémák voltak a rendszerrel. Ezúttal a terheléselosztó és a webszerverek tökéletesen működtek, azonban az adatbázisszerver két rendszerparaméter miatt újraösszeomlott. A szerver újrakonfigurálásával és a rendszer teljes újraindításával a problémát elhárították, és a felhasználók újra megrohanhatták a Neptunt.

A történtek tanulsága, hogy egy ilyen összetett infrastruktúra esetében elég az egyik rétegben bekövetkező hiba, és az az egész rendszerre hatással lehet. Az üzemeltetés automatizált terheléstesztekkel próbálta megkeresni a rendszer szűk keresztmetszetét, és új konfigurációs beállításokkal megszüntetni azokat, de a felhasználói viselkedést nem tudták egyik esetben sem előre meghatározni.

\begin{table}[h]
	\caption{A Neptun rendszer konkurens felhasználóinak száma}
	\centering
	\small
	\begin{tabular}{| p{3.0cm} | p{2.0cm} | p{3.0cm} | p{3.0cm} |}
		\hline
		\rowcolor{MyTableColor} \textbf{Dátum} & \textbf{Jelleg} & \textbf{Max. felhasználó} & \textbf{Op. felhasználó} \\
		\hline
		2010.11.29. 18:00 & vizsga & 7303 & 4623 \\
		\hline
		2010.12.22. 06:00 & tárgy (EO) & 831 & 831 \\
		\hline
		2011.01.10. 18:00 & tárgy & 12062 & 4837 \\
		\hline
		2011.01.12. 18:00 & tárgy (EP) & 1765 & 1765 \\
		\hline
		2011.01.31. 16:00 & tárgy & 1519 & 1519 \\
		\hline
		2011.05.02. 18:00 & vizsga & 2761 & 2761 \\
		\hline
		2011.06.07. 18:00 & tárgy & 6095 & 6095 \\
		\hline
		2011.08.29. 16:00 & tárgy & 1247 & 1247 \\
		\hline
		2011.11.28. 18:00 & vizsga & 4897 & 4897 \\
		\hline
	\end{tabular}
	\normalsize
	\label{tab:neptunusers}
\end{table}


\Aref{tab:neptunusers}.~táblázatban látható, hogy az adott időszakokban hogyan alakult a konkurens felhasználók maximális száma (\textit{Max. felhasználó}) és a rendszer által még kiszolgálható felhasználók száma (\textit{Op. felhasználó}). (Forrás: a BME KTH Neptun Üzemeltetés által kiadott tájékoztatók)


\Aref{fig:neptun_003}.~ábrán látható a 2011. június havi tárgyfelvételi időszak néhány jellemző tárgyfelvétellel kapcsolatos adata. Jól megfigyelhető, hogy az első órában milyen alacsony a sikeres tárgyfelvételek száma, hiszen ebben az időszakban a rendszer nem üzemelt megfelelően, majd a következő oszlopban a már újra konfigurált és újraindított rendszer teljesítménye/felhasználói aktivitása látható.

\begin{figure}[!ht]
\centering
\includegraphics[width=80mm, keepaspectratio]{figures/neptun_003.png}
\caption{Tárgyfelvételek száma a Neptun rendszerben (Forrás: a BME KTH Neptun Üzemeltetés által kiadott tájékoztatók)}
\label{fig:neptun_003}
\end{figure}

A 2011. augusztus végi időszakot \aref{fig:neptun_004}.~ábra mutatja be. Mivel ekkor a rendszer nem volt annyira leterhelve, ezért itt az első oszlop a legmagasabb, majd a másodikban hirtelen csökkenés következett be, végül az idő folyamán fokozatosan csökkent a terhelés. Vagyis a felhasználói kérések száma az első órában volt maximális, amit a rendszer sikeresen kezelt. 

\begin{figure}[!h]
\centering
\includegraphics[width=80mm, keepaspectratio]{figures/neptun_004.png}
\caption{Tárgyfelvételek száma a Neptun rendszerben (Forrás: a BME KTH Neptun Üzemeltetés által kiadott tájékoztatók)}
\label{fig:neptun_004}
\end{figure}

\Aref{fig:neptun_005}.~ábrán a 2011. novemberi tárgyfelvétellel kapcsolatos adatokat láthatjuk. A rendszer kezdeti működési nehézségei itt is jól kivehetők, ha figyelembe vesszük, hogy az egyes oszlopok a sikeres kéréseket mutatják. A rendszer újrakonfigurálásának eredményét a második oszlop mutatja, amely a legmagasabb a vizsgált időszakban.
Szinte mindegyik grafikon esetén megállapítható, hogy a rendszer ez első órában van terhelve, ami a felhasználói viselkedésből ered. A hallgatók tárgyfelvétel esetén szeretnének a legjobb, legnépszerűbb kurzusokra bekerülni, vizsgajelentkezésnél pedig a nekik megfelelő időpontokra jelentkezni, amíg be nem telnek azok.

\begin{figure}[!h]
\centering
\includegraphics[width=80mm, keepaspectratio]{figures/neptun_005.png}
\caption{Tárgyfelvételek száma a Neptun rendszerben (Forrás: a BME KTH Neptun Üzemeltetés által kiadott tájékoztatók)}
\label{fig:neptun_005}
\end{figure}

\subsubsection{Kurzussal kapcsolatos feladatok beadási határideje}

Hallgatók jellemző viselkedése, hogy megpróbálják az utolsó pillanatban elkészíteni a kiadott feladatokat, amelynek következtében a határidő előtti rövid időintervallumban megnövekszik a rendszerben egyszerre tartózkodó felhasználók száma, és állományok feltöltésével esetlegesen túlterhelik a rendszer hálózatát, adatbázisát és nem előrelátó üzemeltetés esetén a tárhely szabad helyének méretében is kritikus szintet érhetnek el.

\subsubsection{Kurzus online teszt, vagy vizsga kitöltés (határ)ideje}

Egy oktatási rendszerben meg lehet határozni olyan időintervallumokat, amelyek egy kurzus tesztjének kitöltésére engedélyezett időszakai. A tesztet az intervallum elejétől a végéig lehetséges kitölteni, azon kívül a teszt kitöltése letiltott funkció. Ezzel a felhasználók arra vannak kényszerítve, hogy az adott intervallumon belül bejelentkezzenek a rendszerbe, és kitöltsék a megadott tesztet.

Mivel az időintervallum meghatározott (az üzemeltetés is könnyen tudomást szerezhet róla), és a várható felhasználói létszám is ismert (az adott kurzuson résztvevő felhasználók száma a maximum), azért viszonylag jól fel lehet készülni erre az erőforrással kapcsolatos igényváltozásra. Természetesen itt sem tudjuk 100\%-osan meghatározni, hogy melyik időpillanatban hány felhasználó lesz a rendszerben, és a felhasználói viselkedéstől is nagyban függ az eloszlás, de maximális értékeket meghatározhatunk. Persze nem biztos, hogy ez költség szempontjából az optimális megoldás.

\subsubsection{Egyéb a kurzussal kapcsolatos offline számonkérés}

Napjainkban egyre több oktatási anyag kerül fel online tárhelyre, amelyet az LMS-ek alapból támogatnak. Egy-egy számonkérés előtt megnövekedhet ezen oktatási anyagok letöltésének száma, amely az anyag méretétől, számától és a kurzusra jelentkezett felhasználók számától függően terhelheti az adatbázist és a hálózatot.

Ebben az esetben a feltöltésekhez hasonlóan nem lehetséges biztos modellt alkotni az erőforrásigények változásáról, hiszen itt is nagyban függünk a felhasználói viselkedéstől. Míg folyamatosan készülő diákok az egész félévben elosztott terhelést jelentenek, addig az utolsó pillanatosok a számonkérés előtt nem sokkal terhelhetik a rendszert. 

\subsubsection{Online előadás közvetítés}

Az Internet térnyerésével és a széles sávú kapcsolatok terjedésével megjelent az \angolul{online stream} lehetősége. Egy ilyen közvetítés egy bizonyos időintervallumban terhelheti a hálózatunkat. Az elérhetőségének korlátozásával (pl. csak az adott kurzust felvett felhasználók nézhetik) a maximális nézőszámot is meg tudjuk előre határozni. Tehát összességében egy ilyen eseményre viszonylag jól fel tudjuk készíteni a rendszerünket erőforrások szempontjából korábbi hálózati kihasználtságok statisztikái alapján.  

\subsubsection{Audiovizuális tananyagokkal rendelkező kurzus számonkérésének ideje}

Ez az eset az online közvetítés és a számonkérések egyvelege. Mivel a számonkérés ideje rögzített, így addig az időpontig számíthatunk arra, hogy a rendszerünket az adott kurzus felhasználói terhelni fogják. A terhelés főleg a hálózatot fogja érinteni.

Nagyon eltérő terheléseket kaphatunk attól függően, hogy az oktatási anyag letölthető-e vagy csak online nézhető. Az előbbi esetben nagyobb az esély az időben jobban szétkenődő erőforrásigény növekedésekre, míg az utóbbinál a határidő előtt nagy konkurens felhasználói számmal és igénynövekedéssel kell számolnunk. Ennek oka lehet, hogy az online elérhető anyagot a számonkérés előtt nem sokkal fogják megnézni a felhasználók, amivel a rendszerünket terhelik. Ellenben a letölthető anyag esetében meg van az a lehetőség, hogy néhány felhasználó már jóval a határidő előtt letölti azt, így csökkentve a maximális felhasználók számát a kritikus időpontban.

Látható, hogy az eltérő oktatási rendszerrel kapcsolatos események különböző módon befolyásolják az egyes erőforrásigényeket. Ezek között vannak olyanok, amelyek csak egy szűkebb felhasználói kört érintenek, ebből kifolyólag a lehetséges maximális igény meghatározható, míg nagyobb csoportok esetén csak körülbelüli értékekkel dolgozhatunk. Ugyanakkor az is megállapítható, hogy a felhasználók viselkedése nagyban befolyásolja a bekövetkező változásokat. 

\section{Modellek}\label{sec:modellek}

Az LMS-ek üzemeltetésének legnagyobb problémája, hogy nem rendelkezünk olyan modellekkel, amelyekkel 100\%-osan előre tudnánk jósolni a rendszer zavartalan működéséhez, az áteresztőképességének szinten tartásához szükséges erőforrások mennyiségét.

A piacon igen sok terhelés tesztelő alkalmazást lehet találni. Ezek közül az egyik legismertebb az Apache JMeter (\href{https://jmeter.apache.org/}{https://jmeter.apache.org/}). Ez egy igen részletesen konfigurálható, elsősorban webes alkalmazások funkcionális terhelés tesztjeire készített asztali alkalmazás.

Valószínűsíthető, hogy a Neptun rendszer üzemeltetése is hasonló szoftverrel próbálta tesztelni az infrastruktúrájának terhelhetőségét, de ahogy láthattuk, ez nem vezetett célra. Ezekkel a szoftverekkel meg tudjuk vizsgálni, hogy a rendszerünk mekkora maximális terhelést visel el, de azt már nem tudjuk meghatározni, hogy a való életben mekkora lesz ez az érték. A probléma itt is, mint az informatika sok más területén, a felhasználó és annak pontosan nem előrejelezhető viselkedése, amelyeket csak statisztikai elemzésekkel tudunk megjósolni.

A Neptun esetében is megfigyelhető, hogy egy kurzus- vagy tárgyfelvételi időszak kezdetén maximális az aktív felhasználók száma, majd ezután az idő elteltével rohamosan csökken. Ez csak általában igaz, hiszen bemutattunk egy olyan esetet is, amikor nem következett be az elvárt roham, mert a felhasználók már korábban elintézhették a tárgyfelvételt, ezért csak az időben szélesen elosztva jelentek meg átlagos értékek a grafikonjainkon, amikor is néhány felhasználó valószínűleg véglegesítette a tárgyait.

Hasonlóan nem teljesen kiszámítható egy kurzus számonkérése előtti felhasználói érdeklődés. Hiszen előfordulhat, hogy két olyan kurzus esetén, amely rendelkezik írásos és audiovizuális oktatási anyaggal, eltérő módon az egyik kurzus esetében csak az írásos segédanyagot, a másik esetben pedig a multimédiás tartalmat fogják igénybe venni, a rendszer szempontjából ezzel eltérő erőforrásigényeket generálva.

\subsection{Lehetséges modell alkotások}

Az oktatási rendszerek esetében egy lehetséges modellalkotási lehetőség lehet az összegyűjtött monitorozási, naplózási adatok statisztikai elemzése. Ezekkel az adatokkal valamennyire jósolható egy esemény által igényelt erőforrások mennyisége, amelynek segítségével akár proaktívan is be lehet avatkozni a rendszer működésébe.

\Aref{fig:neptun_003}., \ref{fig:neptun_004}. és \ref{fig:neptun_005}.~ábrákon látható grafikonokból is le lehet vonni következtetéseket a felhasználók és a rendszer felé irányuló kérések számának eloszlásáról. Hasonló módon a Munin (\href{http://munin-monitoring.org/}{http://munin-monitoring.org/}) nevű grafikus monitorozó eszközzel megoldható, hogy az LMS egyes eseményeit (pl. aktuális felhasználók száma, hálózati eszközök kihasználtsága stb.) grafikonok segítségével jellemezzük. Persze mindezen grafikusan ábrázolt adatoknak csak akkor van értelme, ha mellette azt is tudjuk, hogy az éppen vizsgált időintervallumban milyen oktatással kapcsolatos események történtek (pl. teszt kitöltés, online előadás stb.). 

Tehát az oktatástámogató rendszerekre jellemző erőforrásigény változások esetén a rendszer megfelelő működését úgy tudjuk biztosítani, ha korábbi monitorozási adatok, személyes tapasztalatok, a rendszer használatából eredő események alapján megpróbálunk egy statisztikai elemzésen alapuló modellt felállítani, amelynek segítségével előre tudjuk jelezni az erőforrásigények változását.

\chapter{Információs technológiai infrastruktúrák}
\section{A klasszikus IT infrastruktúra}
A klasszikusnak mondott IT infrastruktúrát az írás korábbi részében már bemutattam. Itt most a háromrétegű architektúra egyes rétegeinek skálázását, megbízhatósági szintjének növelésére adott lehetőségeket szeretném bemutatni.
\subsection{Adatbázis réteg}
\todo{Klaszterek, replikák, stb.}
\subsection{Alkalmazás réteg}
\todo{PHP, Java skálázás?}
\subsection{Webkiszolgáló réteg}
\todo{Loadbalancing} 
\section{Felhőalapú infrastruktúrák}
\subsection{Mi is az a ,,számítási felhő''?}

A ,,számítási felhő'' (angolul \foreignlanguage{english}{cloud computing}) egy modell kényelmes, hálózaton keresztül hozzáférhető, konfigurálható számítási erőforrások (pl. hálózat, szerverek, tárhelyek, alkalmazások és szolgáltatások) egy megosztott készletének elérhetőségére, mely erőforrásokat minimális intézkedési erőfeszítéssel vagy szolgáltatói közbenjárással gyorsan rendelkezésre lehet bocsátani \cite{nistsp800-145}.

\todo{Lehet, hogy még finomítani kell a fordításon!}

\begin{comment}
Cloud computing is a model for enabling ubiquitous, convenient, on-demand network access to a shared pool of configurable computing resources (e.g., networks, servers, storage, applications, and services) that can be rapidly provisioned and released with minimal management effort or service provider interaction. 
\end{comment}

Általában hivatkozás szintjén nincsenek elkülönítve az Interneten keresztül szolgáltatott alkalmazások, és a felhő infrastrukturális részét képező hardverek, szoftverek, amelyek ezeket az alkalmazásokat elérhetővé teszik. Ahogy azt \aref{fig:cloud_computing_hu}.~ábrán is szemléltetni próbálom a felhő részét képezi az alkalmazás, szolgáltatás, és a hardver is.

\begin{figure}[h!]
\centering
\includegraphics[width=1.0\textwidth]{figures/Cloud_computing_hu.png}
\caption{A számítási felhő (\foreignlanguage{english}{cloud computing}) (Forrás: \href{https://en.wikipedia.org/wiki/File:Cloud\_computing.svg}{Wikipedia})} \label{fig:cloud_computing_hu}
\end{figure}

\todo{Object Storage ?= Objektum tároló}

\todo{Queue ?= Üzenetsorok}

\todo{Runtime ?= Számításelosztás}

Mélyebb elemzés során azonban a felhőt rétegekre lehet bontani, amely rétegeket a következő alfejezetben részletezném.

\subsection{A felhő rétegei}

\todo{Ezek lehet inkább szolgáltatás modellek?}

A felhő napjainkban négy architekturális rétegből épül fel, amelyeket alulról fölfelé érdemes megvizsgálni.


\begin{figure}[h!]
\centering
\includegraphics[width=0.25\textwidth]{figures/cloud_retegek.png}
\caption{A számítási felhő rétegei \label{fig:cloud_retegek}}
\end{figure}

\begin{figure}[h!]
\centering
\includegraphics[width=0.5\textwidth]{figures/cloud_service_models.png}
\caption{A számítási felhő szolgáltatási modelljei \label{fig:cloud_service_models}}
\end{figure}
 
\subsubsection{Adattár, mint szolgáltatás (\foreignlanguage{english}{data-Storage-as-a-Service, dSaaS})}
Ezt a réteget \todo{szolgáltatást???} nem minden irodalom szokta említeni, ám én itt mégis külön kezelném, hiszen ez a felhő legalapvetőbb szolgáltatása. Lényege, hogy online tárhelyet biztosít a felhasználóknak. Ilyen szolgáltatást nyújt pl. a \href{http://www.dropbox.com}{Dropbox.com} (főleg személyes felhasználásra, biztonsági mentés, megosztás céljából) vagy az \href{https://aws.amazon.com/s3/}{Amazon S3} (inkább nagy szolgáltatók használják).

A dSaaS oktatási rendszerek esetében sok, nagyméretű adatok esetén lehet előnyös, hiszen nem kell a saját szerverünkön tárolni ezeket, megspórolva ezzel saját adattároló rendszer kialakítását, üzemeltetését. Érdemes megjegyezni, hogy sok adatpéldány esetén is érdemes lehet hasonló szolgáltatás használata, hiszen ebben az esetben az autentikációhoz kötött adatoknál már nem kell a session-öket kezelni, az adatok elérhetőségét megadó URL már tartalmaz egy kódot, csökkentve ezzel a web és alkalmazás szerver terheltségét. 

\todo{dSaaS-ról még valami?}

\todo{Ha a két ábrát össze kellene lőni, ez a legalsó réteg.}

\subsubsection{Infrastuktúra, mint szolgálatás (\foreignlanguage{english}{Infrastructure-as-a-Service, IaaS})}

\todo{IaaS}
\subsubsection{Platform, mint szolgáltatás (\foreignlanguage{english}{Platform-as-a-Service, PaaS})}
\todo{PaaS}
\subsubsection{Szoftver, mint szolgáltatás (\foreignlanguage{english}{Software-as-a-Service,SaaS})}
\todo{SaaS}



\chapter{IT infrastruktúrák proaktív menedzsmentje általános és oktatás támogató rendszerek esetén}

\section{A proaktív menedzsment fogalma}

A rendszer menedzsmentet két típusba sorolhatjuk.  Egy menedzsment rendszert \textit{reaktívnak} mondunk, ha képes gyorsan és hatékonyan reagálni a külső és belső kérelmekre a belső flexibilitás maximalizálásával. Ezt a reaktivitást a rendszer rugalmasságán alapulva decentralizált döntésekkel és a fejlesztés reflexszerű viselkedésével előre definiált szabályok segítségével érik el.\cite{aftsarapms} Tehát egy reaktív menedzsment a rendszerben már bekövetkezett változásokra reagál. A reaktív rendszer vázlata látható \aref{fig:reactive_system}.~ ábrán. A reaktív vezérlés inkább egy cselekvés valamilyen szituációra válaszolva, mint annak a szituációnak a létrehozása, vagy vezérlése.\cite{patariczapro}

\begin{figure}[h!]
\centering
\includegraphics[width=0.80\textwidth]{figures/reactive_system.png}
\caption{A reaktív rendszer áttekintő ábrája (Forrás: \cite{patariczapro}) \label{fig:reactive_system}}
\end{figure}

Egy menedzsment rendszer \textit{proaktív}, ha a reaktív része az előrelátás, illesztés és tanulás folyamataival van kiegészítve, amely folyamatok célja a rendszer támogatása, és annak koherenciájáról és hatékonyságáról való gondoskodás. Egy proaktív rendszer folyamatos monitorozással, előrelátással és tanulással próbál reagálni a rendszerben még be nem következett eseményekre. Egy proaktív rendszer áttekintő vázlatát mutatja be \aref{fig:proactive_system}.~ábra. A poraktív vezérlés inkább egy szituáció irányítása, mint a szituáció által okozott történésekre adott válasz. \cite{patariczapro}

\begin{figure}[h!]
\centering
\includegraphics[width=0.80\textwidth]{figures/proactive_system.png}
\caption{A proaktív rendszer áttekintő ábrája (Forrás: \cite{patariczapro}) \label{fig:proactive_system}}
\end{figure}

\section{Proaktív erőforrás-kezelés általános rendszerek esetén}

\todo{Alakítgatni!}

A proaktív IT rendszerek a reaktívakkal ellentétben előretekintőek, vagyis megpróbálják kitalálni, hogy mi fog történni a rendszerrel a jövőben.

A proaktív üzemeltetés konstans költséggel jár, ezért ott érdemes alkalmazni, ahol a kiesés nagy anyagi kárt okozna.

A reaktív rendszerek reagáló rendszerek, ebből kifolyólag a rendszerben történő problémák nagy költségugrással járnak. 

Egy reaktív rendszer a már bekövetkezett súlyos hibára reagál, míg a proaktív a rendszerben még csak megjelent hibaokra. Ezáltal a hibaok és a végzetes hiba megjelenése között eltelt lappangási időt nyerjünk egy proaktív rendszernél, amely idő alatt pl. át lehet kapcsolni a tartalék erőforrásokra.

Egy proaktív rendszerben állandó megfigyeléseket végzünk, amely megfigyelések során minden releváns esetet figyelembe veszünk. A megfigyelt adatokat analizáljuk, a bemeneti adatokat ellenőrizzük, hogy ténylegesen valós értékeket reprezentálnak-e. Vizuális analízist (pl. párhuzamos koordináták módszere a koherencia megvizsgálására), és automata számítási módszereket alkalmazunk arra, hogy megértsük pontosan mit is jelentenek az adatok. A tudás elemzés során klaszter technikákat, általánosítást használunk, és megpróbálunk előre látni. A szituációt meghatározó fő tényezők tájékoztató jellegűek az akcióterv meghatározásához.

A VMware VMotion technikája a proaktív menedzsment jó példája. Ez a technológiai alkalmazás lehetővé teszi, hogy működés alatt álló virtuális gépeket egyik hardverről egy másikra kerüljenek át a kiesés minimalizálásával. 

\section{Proaktív erőforrás-kezelés oktatás támogató rendszerek esetén}

\subsection{LMS rendszerek proaktív erőforrás-kezelésének automatizálása felhők esetében}

A piacon található nagyobb felhőszolgáltatást nyújtó cégek publikálnak a rendszerükhöz egy API-t (alkalmazás programozási felületet, \angolul{Application Programming Interface}), amely lehetőséget biztosít arra, hogy különféle műveleteket végezzünk el virtuális szerverünkön anélkül, hogy a szolgáltató által nyújtott adminisztrációs felületre belépnénk.

Ilyen műveletek lehetnek az alapvetőnek számító szerver leállítás, elindítás, szerverek kilistázása, vagy a már plusznak számító erőforrások igénylése, módosítása, új terheléselosztó konfigurálása és beüzemelése. A már korábban említett Amazon EC2 szolgáltatás egy viszonylag egyszerű, kényelmes, és funkciókban bővelkedő API-val rendelkezik.

Ebből kifolyólag lehetséges lenne, hogy pl. egy Amazonnál üzemeltetett oktatás támogató rendszer bizonyos eseményekhez előre beállított szabályok alapján újabb erőforrásokat rendeljen hozzá. Például amennyiben az üzemeltetés tudja, hogy egy jövőbeli teszt miatt nagy számú konkurens felhasználóra lehet számítani, akkor az LMS-ben megadott teszt kezdetéhez ütemezve a rendszer automatikusan meghívja azokat az utasításokat, amelynek következtében újabb erőforrásokat kapunk, majd a teszt kitöltésének végén visszaadja azokat az erőforrásokat a szolgáltatónak.

Ennek a módszernek több előnye is van. Egyrészt proaktívan tudjuk kezelni a rendszer terhelésváltozását, megelőzve ezzel a váratlan kieséseket. Másrészt mivel a szolgáltatók nagy része olcsóbban adja az előre igényelt erőforrásokat, mint az igénymegjelenés (\angolul{on-demand}) szerintieket, ezért költségek szempontjából is jobban járunk ezzel a módszerrel.

Az Amazon néhány platformra kész API megvalósításokat nyújt. Ilyen platformok például a PHP, Python, Java, Ruby és .NET, amely platformokat oktatás támogató rendszerek fejlesztésénél is szívesen alkalmazzák, így ezzel is könnyebbséget jelentve a fentebb említett funkciók lefejlesztéséhez.

Az egyedül nehézséget a már korábbi Modellek részben említett problémák adják, vagyis hogy nem egyszerű meghatározni azt, hogy pontosan mennyi erőforrást igényeljünk például egy teszt kitöltésének idejére a szolgáltatótól. 

\subsection{A felhők előnyei és hátrányai LMS-ek üzemeltetése esetén}

Mint ahogy azt már vázoltam egy oktatás támogató rendszert mind hagyományos, mind felhőalapú infrastruktúrán lehet üzemeltetni. Ebben a részben szeretném kicsit összefoglalni, hogy felhő esetében milyen előnyökkel és hátrányokkal jár ez.

Láthattuk, hogy egy LMS erőforrásigénye nagyban függ a felhasználói viselkedéstől, azt is megmutattam, hogy napjainkban nincsenek még 100\%-os modellek, módszerek ezek előrejelzésére. Amíg a klasszikus IT infrastruktúránál az üzemeltetés csak próbálja mindig meghatározni a maximális erőforrásigényt, és ennek megfelelően kialakítani, fejleszteni a rendszert, addig a felhők proaktívitásának köszönhetően ezek a plusz erőforrásigények automatikusan kielégítésre kerülnek, igaz, hogy költség szempontjából kicsit drágábbak. Viszont ezeknek a költségeknek is több összetevője van, hiszen a felhő alapú üzemeltetés esetén nem kell foglalkoznunk eszközbeszerzéssel, karbantartással, hibás hardverek pótlásával, amely a nem várt költségek nagy részét tehetik ki. Felhők esetén előre ki tudjuk számolni a költségeinket, a rendszer megfelelő méretezésével pedig a drágább \angolul {on-demand} erőforrások költségeit is minimalizálhatjuk.

A felhők hátrányaként jelenik az saját IT infrastruktúra üzemeltetésével szemben, hogy minden adat a mi tulajdonunkban van, tehát biztonságuk csak tőlünk függ. Előfordulhatnak olyan adatok (pl. személyes adatok, vagy kutatások eredményei), amelyeket nem szeretnénk kiadni harmadik félnek. Ilyen esetekben egy publikus felhőszolgáltatást nyújtó szolgáltatónál nem üzemeltethetjük a rendszerünket.

Természetesen találunk olyan problémákat is, amelyek mindkét struktúra típus esetén megtalálható, mint például előfordulhatnak olyan hálózati problémák, amelyek miatt a rendszerünk nem lesz elérhető se a klasszikus esetben (pl. egyetemi router kiesése), se a felhőalapúban (a szolgáltatónk telephelyén kiesik egy fontos hálózati elem).

\chapter{Összefoglalás}

Írásomban megpróbáltam megvizsgálni, hogy az oktatástámogató rendszerek esetében mennyire lehet alkalmazni az általános rendszerek esetén alkalmazott rendszerfelügyelettel kapcsolatos technikákat az erőforrásigények változása szempontjából.

Először bemutattam, hogy mik azok a célalkalmazások, amelyeket vizsgálok, hol alkalmazzák őket, milyen technológiákra épülnek, majd bemutattam, hogy egy informatikai infrastruktúra esetében mik az oktatástámogató rendszerek által támasztott igények. Leírtam, hogy egy példa rendszer, a Moodle LMS milyen felépítéssel és igényekkel rendelkezik.

A következő fejezetben összegyűjtöttem, hogy az oktatástámogató rendszerek milyen erőforrásigényekkel rendelkeznek, hogyan változhatnak ezek az igények, és mik lehetnek ezen változások okozói. A fejezet végén a lehetséges modelleket próbáltam meg leírni.

Ezután a klasszikus IT infrastruktúra nagy rendelkezésre állás elérésének lehetőségeit szedtem össze, majd a fejezet nagy részében foglalkoztam a napjainkban igen elterjedt felhőalapú infrastruktúrával és annak szolgáltatásaival. Igyekeztem az LMS-ek szempontjából vizsgálni az egyes részeket.

Szakdolgozatom utolsó részében arra kerestem a választ, hogy miben különbözik egy reaktív és proaktív infrastruktúra menedzsment, és az oktatástámogató rendszerek üzemeltetését hogyan segítheti a proaktív megoldás.

\todo{Konklúzió? Folytatási lehetőségek?}

%----------------------------------------------------------------------------
\chapter*{Köszönetnyilvánítás}\addcontentsline{toc}{chapter}{Köszönetnyilvánítás}
%----------------------------------------------------------------------------

Azért, hogy ez az írás elkészült, köszönettel tartozom konzulensemnek Huszerl Gábornak, aki segített a szakdolgozat témájának, felépítésének kitalálásában. Kocsis Imrének a szerverek erőforrásigényeivel foglalkozó és Gönczy Lászlónak a proaktív menedzsmenttel foglalkozó forrásanyagért. Dr. Pataricza Andrásnak a lelkesítő szavakért.

Testvéreimnek Nádudvari Zoltánnak, akivel együtt tartottuk egymásban a lelket a saját szakdolgozatunk írása közben, Nádudvari Tamásnak, aki segített megbirkózni a szövegek fordításával. Végezetül szeretnék köszönetet mondani szüleimnek, és néhány jó barátnak, akik többek között kitartásra buzdítottak a félév során. 

\listoffigures\addcontentsline{toc}{chapter}{Ábrák jegyzéke}
\listoftables\addcontentsline{toc}{chapter}{Táblázatok jegyzéke}

\bibliography{irodalomjegyzek}
\addcontentsline{toc}{chapter}{Irodalomjegyzék}
\bibliographystyle{plain}

%----------------------------------------------------------------------------
\appendix
%----------------------------------------------------------------------------
\chapter*{Függelék}\addcontentsline{toc}{chapter}{Függelék}
\setcounter{chapter}{6}  % a fofejezet-szamlalo az angol ABC 6. betuje (F) lesz
\setcounter{equation}{0} % a fofejezet-szamlalo az angol ABC 6. betuje (F) lesz
\numberwithin{equation}{section}
\numberwithin{figure}{section}
\numberwithin{lstlisting}{section}
%\numberwithin{tabular}{section}

%----------------------------------------------------------------------------
\section{Az Amazon AWS néhány fontosabb API metódusa}\label{sec:appendix_amazonaws}
%----------------------------------------------------------------------------

A következő táblázatokban az Amazon Web Services (AWS) szolgáltatásai nyújtott API egyik megvalósításának néhány metódusát gyűjtöttem össze. A megvalósítás a Python alapú Boto (forráskód: \href{https://github.com/boto/boto}{https://github.com/boto/boto}, dokumentáció: \href{http://readthedocs.org/docs/boto/en/latest/}{http://readthedocs.org/docs/boto/en/latest/}) nevű projekt.

\subsection{Amazon EC2 API metódusok}

\begin{table}[h]
	\caption{Az Amazon EC2 API boto.ec2.connection osztályának néhány metódusa (associate\_address - disassociate\_address)}
	\centering
	\small
	\begin{tabular}{| p{7.5cm} | p{6.5cm} |}
		\hline
		\rowcolor{MyTableColor} \textbf{Metódus/függvény} & \textbf{Leírás} \\
		\hline
		\texttt{\textbf{associate\_address}(instance\_id, public\_ip=None, allocation\_id=None)} & Egy IP címet (Elastic IP-t) rendel az éppen futó szerver példányunkhoz. \\
		\hline
		\texttt{\textbf{attach\_volume}(volume\_id, instance\_id, device)} & A \texttt{volume\_id} EBS (Elastic Block Store) kötetet az \texttt{instance\_id} szerverpéldányhoz csatolja. \\
		\hline
		\texttt{\textbf{create\_image}(instance\_id, name, description=None, no\_reboot=False)} & Egy AMI-t (Amazon Machine Image) készít az \texttt{instance\_id} szerverpéldány aktuális futó vagy leállított állapotából. \\
		\hline
    	\texttt{\textbf{create\_key\_pair}(key\_name)} & Egy új kulcspárt készít, amely kulcs az azonosításhoz szükséges. \\
		\hline
		\texttt{\textbf{create\_snapshot}(volume\_id, description=None)} & Egy létező EBS kötetről készít pillanatképet. \\
		\hline
		\texttt{\textbf{create\_tags}(resource\_ids, tags)} & Új metaadat címkéket készít a \texttt{resource\_ids} azonosítójú erőforrásokhoz. \\
		\hline
        \texttt{\textbf{create\_volume}(size, zone, snapshot=None)} & Egy új \texttt{size} méretű EBS kötetet készít. \\
        \hline
        \texttt{\textbf{delete\_key\_pair}(key\_name)} & Törli a \texttt{key\_name} nevű azonosító kulcspárt. \\
        \hline
        \texttt{\textbf{delete\_snapshot}(snapshot\_id)} & Törli a \texttt{snapshot\_id} azonosítójú pillanatképet. \\
        \hline        
        \texttt{\textbf{delete\_tags}(resource\_ids, tags)} & Tölri a metaadat címkéket a \texttt{resource\_ids} azonosító erőforrásokról. \\
        \hline
        \texttt{\textbf{delete\_volume}(volume\_id)} & Töröli a \texttt{volume\_id} azonosítójú EBS kötetet. \\
        \hline
        \texttt{\textbf{detach\_volume}(volume\_id, instance\_id=None, device=None, force=False)} & Lecsatol egy EBS kötetet egy EC2 példányról. \\
        \hline
        \texttt{\textbf{disassociate\_address}(public\_ip=None, association\_id=None)} & Elveszi az Elastic IP címet az aktuálisan futó példánytól. \\        
        \hline
    \end{tabular}
	\normalsize
	\label{tab:aws_boto_ec2_01}
\end{table}        
        
\begin{table}[h]
	\caption{Az Amazon EC2 API boto.ec2.connection osztályának néhány metódusa (get\_all\_images - reset\_snapshot\_attribute)}
	\centering
	\small
	\begin{tabular}{| p{8.0cm} | p{6.0cm} |}
	\hline
		\rowcolor{MyTableColor} \textbf{Metódus/függvény} & \textbf{Leírás} \\
	    \hline          
        \texttt{\textbf{get\_all\_images}(image\_ids=None, owners=None, executable\_by=None, filters=None)} & Lekéri az összes EC2 szerver image-t, amelyek a fiókhoz kapcsolódnak. \\
        \hline   
        \texttt{\textbf{get\_all\_instances}(instance\_ids=None, filters=None)} & Lekéri az összes szerver példányt, amelyek a fiókhoz tartoznak. \\
        \hline
        \texttt{\textbf{get\_all\_key\_pairs}(keynames=None, filters=None)} & Visszaadja a fiókhoz tartozó összes kulcspárt. \\
        \hline
        \texttt{\textbf{get\_all\_ramdisks}(ramdisk\_ids=None, owners=None)} & Lekéri az összes EC2 ramdisket,a melyek elérhetőek a fiókból. \\
        \hline
        \texttt{\textbf{get\_all\_regions}(region\_names=None, filters=None)} & Lekéri az összes elérhető régióját az EC2 szolgáltatásnak. \\
        \hline
        \texttt{\textbf{get\_all\_volumes}(volume\_ids=None, filters=None)} & Lekéri az összes kötetet. \\
        \hline
        \texttt{\textbf{get\_console\_output}(instance\_id)} & Letölti az adott szerver példány konzol kimenetét. \\
        \hline
        \texttt{\textbf{get\_image\_attribute}(image\_id, attribute='launchPermission')} & Lekéri egy szerver képállomány attribútumait. \\
        \hline
        \texttt{\textbf{get\_instance\_attribute}(instance\_id, attribute)} & Lekéri egy példány attribútumait. \\
        \hline
        \texttt{\textbf{get\_password\_data}(instance\_id)} & Egy Windows-os szerverpéldány titkosított jelszavát kéri le. \\
        \hline
        \texttt{\textbf{get\_snapshot\_attribute}(snapshot\_id, attribute='createVolumePermission')} & Lekéri egy pillanatkép attribútumának információját. \\
        \hline
        \texttt{\textbf{import\_key\_pair}(key\_name, public\_key\_material)} & Egy külső program által generált RSA kulcspár publikus kulcsát importálja be. \\
        \hline
        \texttt{\textbf{modify\_image\_attribute}(image\_id, attribute='launchPermission', operation='add', user\_ids=None, groups=None, product\_codes=None)} & Egy szerver képállomány attribútumait változtatja meg. \\
        \hline
        \texttt{\textbf{modify\_instance\_attribute}(instance\_id, attribute, value)} & Egy szerverpéldány attribútumait változtatja meg. \\
        \hline
        \texttt{\textbf{modify\_snapshot\_attribute}(snapshot\_id, attribute='createVolumePermission', operation='add', user\_ids=None, groups=None)} & Egy pillanatkép attribútumát változtatja meg. \\
        \hline
        \texttt{\textbf{monitor\_instances}(instance\_ids)} & Engedélyezi a CloudWatch monitorozót a megadott szerver példányokon. \\
        \hline
        \texttt{\textbf{reboot\_instances}(instance\_ids=None)} & Újraindítja a megadott példányokat. \\
        \hline
        \texttt{\textbf{release\_address}(public\_ip=None, allocation\_id=None)} & Felszabadít egy Elastic IP-címet. \\
        \hline
        \texttt{\textbf{reset\_image\_attribute}(image\_id, attribute='launchPermission')} & Visszaállítja egy AMI attribútumának alapértelmezett értékét. \\
        \hline
        \texttt{\textbf{reset\_instance\_attribute}(instance\_id, attribute)} & Visszaállítja egy példány attribútumának alapértelmezett értékét. \\
        \hline
        \texttt{\textbf{reset\_snapshot\_attribute}(snapshot\_id, attribute='createVolumePermission')} & Visszaállítja egy pillanatkép attribútumának alapértelmezett értékét. \\
        \hline
    \end{tabular}
	\normalsize
	\label{tab:aws_boto_ec2_02}
\end{table}        
        
\begin{table}[h]
	\caption{Az Amazon EC2 API boto.ec2.connection osztályának néhány metódusa (run\_instances - unmonitor\_instances)}
	\centering
	\small
	\begin{tabular}{| p{9.0cm} | p{5.0cm} |}
	    \hline
        \rowcolor{MyTableColor} \textbf{Metódus/függvény} & \textbf{Leírás} \\
	    \hline 
        \texttt{\textbf{run\_instances}(image\_id, min\_count=1, max\_count=1, key\_name=None, security\_groups=None, user\_data=None, addressing\_type=None, instance\_type='m1.small', placement=None, kernel\_id=None, ramdisk\_id=None, monitoring\_enabled=False, subnet\_id=None, block\_device\_map=None, disable\_api\_termination=False, instance\_initiated\_shutdown\_behavior=None, private\_ip\_address=None, placement\_group=None, client\_token=None, security\_group\_ids=None)} & Egy szerver képállományt (image-t) futtat az EC2-n. \\
        \hline
        \texttt{\textbf{start\_instances}(instance\_ids=None)} & Elindítja az adott példányt. \\
        \hline
        \texttt{\textbf{stop\_instances}(instance\_ids=None, force=False)} & Leállítja az adott példányt. \\
        \hline
        \texttt{\textbf{terminate\_instances}(instance\_ids=None)} & Megszünteti az adott példány futását. \\
        \hline
        \texttt{\textbf{unmonitor\_instances}(instance\_ids)} & Kikapcsolja a CloudWatch monitorozót az adott példányokon. \\
        \hline
	\end{tabular}
	\normalsize
	\label{tab:aws_boto_ec2_03}
\end{table}

\label{page:last}
\end{document}
